\chapter{混杂}

设想一位研究者做了一项观察性研究，想回答一个因果问题：“一个人抬头望天，是否会让别的行人也抬头望天？”她发现第一个行人抬头与第二个行人抬头之间存在关联。但她也发现，行人往往是在听到头顶一声雷鸣时才抬头。于是，究竟是什么让第二个行人抬头——是第一个行人抬头，还是那声雷鸣——并不清楚。她得出结论：一个人抬头所产生的效应，被一声雷鸣的存在所混杂。

在随机化实验中，处理由掷硬币决定；而在观察性研究中，处理（例如一个人抬头）可能由许多因素（例如一声雷鸣）决定。若这些因素会影响结局发生的风险（例如另一个行人抬头），那么这些因素的效应便与处理的效应纠缠在一起。这时我们说存在混杂，它不过是受处理者与未受处理者之间缺乏可交换性的一种形式。混杂常被视为观察性研究的主要缺陷。存在混杂时，即便研究总体任意大，“关联不是因果”这句老话依然成立。本章给出混杂的定义，并回顾校正混杂的方法。

\section{混杂的结构}

混杂的结构，即由处理与结局的共同原因所致的偏倚，可以用因果图表示。例如，图~7.1（与图~6.1 相同）描绘了处理 $A$、结局 $Y$，以及二者的共享（共同）原因 $L$。该图显示处理与结局之间关联的两个来源：1）路径 $A\to Y$，它代表 $A$ 对 $Y$ 的因果效应；2）$A$ 与 $Y$ 之间包含共同原因 $L$ 的路径 $A\leftarrow L\to Y$。通过共同原因 $L$ 连接 $A$ 与 $Y$ 的路径 $A\leftarrow L\to Y$，是后门路径的一个例子。\mnote{在因果 DAG 中，后门路径是处理与结局之间的一条非因果路径；即使删去所有从处理指向其他变量（处理的子孙）的箭头，这条路径依然存在。也就是说，路径上有一个箭头指向处理。}

\begin{figure}[htbp]
\centering
\begin{tikzpicture}[causal,shorten >=3pt,shorten <=1.5pt,every node/.style={font=\small}]
  \node (L) at (0,0) {$L$};
  \node (A) at (1.6,0) {$A$};
  \node (Y) at (3.2,0) {$Y$};
  \draw[ed] (L) -- (A);
  \draw[ed] (A) -- (Y);
  \draw[ed] (L) to[out=55,in=125] (Y);
\end{tikzpicture}
\caption{处理与结局的共同原因}
\label{fig:7.1}
\end{figure}

如果图~7.1 中不存在共同原因 $L$，那么处理与结局之间唯一的路径就是 $A\to Y$，于是 $A$ 与 $Y$ 之间的全部关联都归因于 $A$ 对 $Y$ 的因果效应。也就是说，关联风险比 $\Pr[Y=1\given A=1]/\Pr[Y=1\given A=0]$ 将等于因果风险比 $\Pr[Y^{a=1}=1]/\Pr[Y^{a=0}=1]$；关联就是因果。但共同原因 $L$ 的存在，为处理 $A$ 与结局 $Y$ 之间制造了额外的关联来源，我们称之为 $A$ 对 $Y$ 之效应的混杂。由于混杂，关联风险比不等于因果风险比；关联不是因果。

观察性研究中混杂的例子比比皆是。考虑以下几种处理对健康结局之效应的混杂：

\begin{itemize}
\item 职业因素：若“身体健康”$L$ 既是积极消防员的原因，又是较低死亡风险的原因，则“消防员”这一职业 $A$ 对死亡风险 $Y$ 的效应会受到混杂。这种偏倚（图~7.1 的因果图所描绘）常被称为健康工人偏倚。
\item 临床决策：若药物 $A$（比如阿司匹林）更可能开给具有某种状况 $L$（比如心脏病）的个体，而 $L$ 既是该治疗的适应证，又是该病的风险因素，则药物 $A$ 对疾病 $Y$（比如卒中）风险的效应会受到混杂。心脏病 $L$ 是卒中 $Y$ 的风险因素，因为 $L$ 对 $Y$ 有直接因果效应（如图~7.1），或者如图~7.2 所示，因为 $L$ 与 $Y$ 都由一个未测量的变量——动脉粥样硬化 $U$——所致。这种偏倚称为适应证混杂或渠道偏倚；后一术语常专用于描述由病人特有的风险因素 $L$ 造成的偏倚——这些因素促使医生在某一类药中选用特定的药物 $A$。
\item 生活方式：若行为 $A$（比如锻炼）与另一种行为 $L$（比如吸烟）相关，而 $L$ 对 $Y$ 有因果效应，且 $L$ 往往与 $A$ 同时出现，则行为 $A$ 对 $Y$（比如死亡）风险的效应会受到混杂。变量 $L$、$A$、$Y$ 的结构如图~7.3 的因果图所示，其中未测量的变量 $U$ 代表导致缺乏锻炼与吸烟的某种性格与社会因素。另一个常见问题：亚临床疾病 $U$ 同时导致缺乏锻炼 $A$ 与临床疾病 $Y$ 的风险升高。当 $L$ 未知时，这种形式的混杂常被称为反向因果。
\item 遗传因素：若存在一段 DNA 序列 $L$，它对 $Y$ 有因果效应，且在携带 $A$ 的人中更常见，则 DNA 序列 $A$ 对某种性状 $Y$ 之风险的效应会受到混杂。这种偏倚同样由图~7.3 的因果图表示，称为连锁不平衡或群体分层；后一术语常专用于描述在来自不同族群个体混合而成的样本中做研究时产生的偏倚。因此，变量 $U$ 可以代表族群或其他导致 DNA 序列连锁的因素。
\item 社会因素：若 55 岁时的失能程度 $L$ 同时影响未来的收入与 75 岁时的失能水平，则 65 岁时的收入 $A$ 对 75 岁时失能水平 $Y$ 的效应会受到混杂。这种偏倚可用图~7.1 的因果图描绘。
\item 环境暴露：若其他污染物 $L$ 的水平与 $A$ 的水平共同变化，且 $L$ 导致冠心病，则空气中颗粒物 $A$ 对冠心病 $Y$ 之风险的效应会受到混杂。这种偏倚同样由图~7.3 的因果图表示，其中未测量的变量 $U$ 代表影响各类空气污染水平的天气状况。
\end{itemize}

\begin{figure}[htbp]
\centering
\begin{tikzpicture}[causal,shorten >=3pt,shorten <=1.5pt,every node/.style={font=\small}]
  \node (U) at (0,-1.3) {$U$};
  \node (L) at (0,0) {$L$};
  \node (A) at (1.6,0) {$A$};
  \node (Y) at (3.2,0) {$Y$};
  \draw[ed] (U) -- (L);
  \draw[ed] (U) -- (Y);
  \draw[ed] (L) -- (A);
  \draw[ed] (A) -- (Y);
\end{tikzpicture}
\caption{未测量的共同原因}
\label{fig:7.2}
\end{figure}

\mnote{有些作者倾向于把未测量的共同原因 $U$（以及从它出发的两个箭头）替换为它所导致的那些已测变量之间的一条双向边。}

\begin{figure}[htbp]
\centering
\begin{tikzpicture}[causal,shorten >=3pt,shorten <=1.5pt,every node/.style={font=\small}]
  \node (U) at (0,-1.5) {$U$};
  \node (L) at (0,0) {$L$};
  \node (A) at (1.6,0) {$A$};
  \node (Y) at (3.2,0) {$Y$};
  \draw[ed] (U) -- (L);
  \draw[ed] (U) -- (A);
  \draw[ed] (A) -- (Y);
  \draw[ed] (L) to[out=55,in=125] (Y);
\end{tikzpicture}
\caption{共享未测量原因的两个已测原因}
\label{fig:7.3}
\end{figure}

在所有这些情形中，偏倚具有相同的结构：它源于存在一个同时为处理 $A$ 与结局 $Y$ 所共享的原因（$L$ 或 $U$），从而在 $A$ 与 $Y$ 之间留下一条开放的后门路径。我们把由处理与结局的共同原因所致的偏倚称为混杂，而用其他名称指称由处理与结局共同原因之外的结构性原因所致的偏倚。为表述简便，本章通篇假定正值性与一致性成立，因果图中的所有节点都被完美测量，不存在带方框的选择节点 $S$（即数据是所关心总体的随机样本），且不存在随机变异。带选择节点的因果图将在第 8 章讨论，带测量误差节点的因果图将在第 9 章讨论。随机变异将在第 10 章讨论。\mnote{混杂的早期统计描述，离散变量由 Yule (1903) 给出，连续变量由 Pearson 等 (1899) 给出。Yule 把由混杂造成的关联称为“虚构的”“虚幻的”与“表面上的”；Pearson 等 (1899) 称之为“虚假”相关。然而，这些关联一点也不虚构、虚幻、表面或虚假。由共同原因造成的关联是实实在在的关联，尽管不能把它们因果地解释为处理效应。或者用 Yule 的话说，它们是“不能赋予最显而易见的物理意义”的关联。}

\section{混杂与可交换性}

现在我们把用因果图定义的混杂概念，与此前各章用反事实定义的可交换性概念联系起来。

当可交换性 $Y^{a}\indep A$ 成立时（如在全个体接受治疗概率相同的边际随机化实验中），无需对任何变量做校正即可识别平均因果效应。对二值处理 $A$，平均因果效应 $\E[Y^{a=1}]-\E[Y^{a=0}]$ 等于条件均值之差 $\E[Y\given A=1]-\E[Y\given A=0]$。当可交换性 $Y^{a}\indep A$ 不成立、但条件可交换性 $Y^{a}\indep A\given L$ 成立时（如在接受治疗的概率随 $L$ 取值而变的条件随机化实验中），平均因果效应同样可以识别。然而，正如第 2 章所述，要在总体中识别因果效应 $\E[Y^{a=1}]-\E[Y^{a=0}]$，需要通过标准化或 IP 加权对变量 $L$ 做校正。此外，正如第 4 章所述，条件可交换性还允许通过分层识别任何取值 $l$ 下的条件因果效应 $\E[Y^{a=1}\given L=l]-\E[Y^{a=0}\given L=l]$。\marginnote{\textcolor{black!85}{关于混杂与可交换性之间关系的详细讨论，见 Greenland and Robins (1986, 2009)。}}[1.2cm]

\marginnote{\makebox[0pt][l]{\parbox[t]{3.3cm}{\raggedright\scriptsize 在条件可交换性下，$\E[Y^{a}]=$\\ $\sum_l \E[Y\given A=a,L=l]$\\ ${}\times\Pr[L=l]$。}}}[2.6cm]

在实践中，若我们相信混杂可能存在，一个关键问题便出现了：能否确定是否存在一组已测量的协变量 $L$，使条件可交换性成立？回答这个问题很难，因为在复杂的因果系统中，用条件可交换性 $Y^{a}\indep A\given L$ 来思考往往并不直观。

本章将说明，如果知道产生数据的因果 DAG，回答这个问题是可能的。为此，假定我们知道真实的因果 DAG（眼下它从何而来并不重要：或许我们有充分的领域知识，或许有一位全知的神把它交给了我们）。因果 DAG 如何让我们判定是否存在一组变量 $L$ 使条件可交换性成立？主要有两种途径：(i) 对因果 DAG 应用后门准则；(ii) 把因果 DAG 变换为 SWIG。使用 SWIG 是更直接的途径，但它也需要更多工具，所以我们先解释后门准则；SWIG 途径将在 7.5 节描述。\mnote{Pearl (1995, 2009) 提出了用于因果效应非参数识别的后门准则。}

若一组协变量 $L$ 阻断了 $A$ 与 $Y$ 之间所有后门路径，且 $L$ 不含任何为处理 $A$ 之子孙的变量，则称 $L$ 满足后门准则。在忠实性以及技术注记~7.1 所讨论的另一条件之下，条件可交换性 $Y^{a}\indep A\given L$ 成立当且仅当 $L$ 满足后门准则。（下文将基于 SWIG 给出这一事实的简单证明。）因此，对于给定的一组协变量 $L$，条件可交换性是否成立，我们现在都能回答。于是，通过尝试处理的全部已测非子孙变量的每一个子集，我们可以回答任何子集是否满足条件可交换性的问题。（事实上，已有算法能大幅减少需要尝试的子集数目。）

\renewcommand{\TPname}{技术注记（Technical Point）}
\begin{techpoint}[条件可交换性是否蕴含后门准则？]
$L$ 满足后门准则，总是蕴含给定 $L$ 的条件可交换性，即便没有忠实性。正文中我们还说，在忠实性之下，给定 $L$ 的条件可交换性蕴含 $L$ 满足后门准则。最后这句话在 FFRCISTG 模型下成立（见技术注记~6.2）。相反，在 NPSEM-IE 模型下，即使后门准则不成立，条件可交换性也可能成立，例如节点为 $A$、$L$、$Y$ 且箭头为 $A\to L$、$A\to Y$ 的因果 DAG。本书除非另有说明，总是假定 FFRCISTG 模型与忠实性。

因果模型之间的这一差别，源于 NPSEM-IE 与 FFRCISTG 模型不同，它假定了反事实之间的跨世界独立性。然而，跨世界独立性永远无法被任何随机化实验验证，哪怕在原则上也不行；这正是 Robins (1986, 1987) 在其 FFRCISTG 模型中不假定跨世界独立性的原因。我们将在第 23 章回到这一问题。
\end{techpoint}
\renewcommand{\TPname}{技术注记}

现在把后门准则（即可交换性）与混杂联系起来。后门准则得到满足的情形有两种：

\begin{enumerate}
\item 处理与结局没有共同原因。在图~6.2 中，处理与结局没有共同原因，因而没有需要阻断的后门路径。此时满足后门准则的变量集是空集，我们说没有混杂。
\item 处理与结局有共同原因，但 $A$ 的已测非子孙变量的一个子集 $L$ 足以阻断所有后门路径。在图~7.1 中，满足后门准则的变量集是 $L$。于是我们说存在混杂，但不存在残余混杂——消除它需要对未测量变量做校正（而这当然是不可能的）。为简洁起见，我们说没有未测量混杂。
\end{enumerate}

第一种情形描述的是边际随机化实验：其中不应期望存在混杂，因为治疗分配完全由掷硬币——或其计算机化升级版：随机数发生器——决定，而掷硬币不可能导致结局。也就是说，当处理被无条件随机分配时，受处理者与未受处理者预期是可交换的，因为不存在共同原因，等价地，因为不存在开放的后门路径；边际可交换性，即 $Y^{a}\indep A$，等价于处理与结局没有共同原因。

第二种情形描述的是条件随机化实验：其中具有相同 $L$ 取值的所有个体接受治疗的概率相同，但按设计，这一概率随 $L$ 的取值而变，即存在箭头 $L\to A$。若 $L$ 同时是结局的原因（如图~7.1），或是结局之某个未测量原因的子孙（如图~7.2），那么这种实验设计就保证了混杂的存在，于是存在开放的后门路径。然而，对协变量 $L$ 做条件化将阻断所有后门路径，因而条件可交换性，即 $Y^{a}\indep A\given L$，将会成立。当对 $L$ 做条件化能阻断所有后门路径——即在 $L$ 的各个水平内受处理者与未受处理者是可交换的——我们说 $A$ 的一个已测非子孙变量集 $L$ 是混杂校正的充分集。

以我们的心脏移植研究——一项条件随机化实验——为例。接受移植（$A=1$）的个体与其他人（$A=0$）不同，因为如果受治疗者未受治疗，其死亡风险 $Y$ 本会高于那些实际未受治疗者：受治疗者中严重心脏病 $L$ 的比例更高，而 $L$ 是 $A$ 与 $Y$ 的共同原因。处理与结局存在共同原因，意味着受处理者与未处理者不是边际可交换的，但在给定 $L$ 时是条件可交换的。第二种情形也正是人们在测量了许多变量 $L$ 的观察性研究中所期望的。

后门准则并不回答关于混杂大小或方向的问题。从逻辑上讲，可能有某些未被阻断的后门路径很弱（例如 $L$ 对 $A$ 或 $Y$ 的效应都不大），因而只引起很小的偏倚；也可能有若干条很强的后门路径把偏倚推向相反方向，从而净偏倚很弱。由于未测量混杂不是“全有或全无”的问题，在实践中考虑偏倚的预期方向与大小很重要（见精要~7.1）。

\section{混杂与后门准则}

下面我们举几个应用后门准则的例子，以判定 $A$ 对 $Y$ 的因果效应是否可识别；若可识别，又需要哪些变量来保证条件可交换性。请记住，本章所有因果 DAG 中的节点都被完美测量，且未对其做条件化。

图~7.1 中存在混杂，因为处理 $A$ 与结局 $Y$ 共享原因 $L$，即 $A$ 与 $Y$ 之间经由 $L$ 存在一条开放的后门路径。然而，对 $L$ 做条件化可以阻断这条后门路径。因此，如果研究者收集了所有个体关于 $L$ 的数据，则在给定 $L$ 时不存在未测量混杂。

图~7.2 中存在混杂，因为处理 $A$ 与结局 $Y$ 共享未测量的原因 $U$，即 $A$ 与 $Y$ 之间经由 $U$ 存在一条后门路径。（与变量 $L$、$A$、$Y$ 不同，变量 $U$ 未被研究者测量。）如果收集了 $U$ 的数据，理论上可以通过对 $U$ 做条件化阻断这条后门路径，从而消除混杂。然而，对 $L$ 做条件化同样可以阻断这条后门路径。因此，给定 $L$ 时不存在未测量混杂。

图~7.3 中同样存在混杂，因为处理 $A$ 与结局 $Y$ 共享原因 $U$；而后门路径同样可以通过对 $L$ 做条件化来阻断。因此，给定 $L$ 时不存在未测量混杂。

现在考虑图~7.4。该因果图中不存在处理 $A$ 与结局 $Y$ 的共同原因，因而没有混杂。$A$ 与 $Y$ 之间经由 $L$ 的后门路径（$A\leftarrow U_2\to L\leftarrow U_1\to Y$）被阻断，因为 $L$ 在这条路径上是一个对撞节点。于是 $A$ 与 $Y$ 之间的全部关联都归因于 $A$ 对 $Y$ 的效应：关联就是因果。例如，设 $A$ 代表体力活动，$Y$ 代表宫颈癌，$U_1$ 代表癌前病变，$L$ 代表针对癌前病变的诊断检查（巴氏涂片），$U_2$ 代表注重健康的性格（更爱运动，也更常看医生）。那么在图~7.4 的因果图下，体力活动 $A$ 对癌症 $Y$ 的效应无混杂，计算 $\Pr[Y^{a=1}]$ 或 $\Pr[Y^{a=0}]$、进而计算总体中的因果效应，都不需要对 $L$ 做校正。

\begin{figure}[htbp]
\centering
\begin{tikzpicture}[causal,shorten >=3pt,shorten <=1.5pt,every node/.style={font=\small}]
  \node (U1) at (0,1.2) {$U_1$};
  \node (L)  at (0,0) {$L$};
  \node (U2) at (0,-1.2) {$U_2$};
  \node (A)  at (1.7,0) {$A$};
  \node (Y)  at (3.4,0) {$Y$};
  \draw[ed] (U1) -- (L);
  \draw[ed] (U1) -- (Y);
  \draw[ed] (U2) -- (L);
  \draw[ed] (U2) -- (A);
  \draw[ed] (A) -- (Y);
\end{tikzpicture}
\caption{没有共同原因，$L$ 是路径上的对撞节点}
\label{fig:7.4}
\end{figure}

\renewcommand{\FPname}{精要（Fine Point）}
\begin{finepoint}[混杂偏倚的强度与方向]
设想你做了一项观察性研究，以识别心脏移植 $A$ 对死亡 $Y$ 的效应，并假定不存在未测量混杂。一位思虑周密的批评者说：“这项观察性研究的推断可能因吸烟 $L$ 的潜在混杂而错误。”一个关键问题是：这一偏倚会使心脏移植效应的估计被削弱还是被夸大。例如，设你的研究得到的风险比为 0.6（据估计，心脏移植使随访期死亡率降低 40\%），并且正如审稿人所怀疑的，吸烟 $L$ 是 $A$（吸烟者较少接受心脏移植）与 $Y$（吸烟者更可能死亡）的共同原因。由于心脏移植组（$A=1$）中吸烟者（$L=1$）少于另一组（$A=0$），即便在“治疗 $A$ 对 $Y$ 无效应”的零假设下，人们也会预期在 $A=1$ 组中观察到更低的死亡风险。因此，对吸烟做校正会使效应估计上移（比如从 0.6 移到 0.7）。换言之，未对吸烟做校正，导致心脏移植的有益平均因果效应被夸大。

预测混杂偏倚方向的一种方法是使用带符号的因果图。考虑图~7.1 的因果图，$L$、$A$、$Y$ 均为二值变量。若 $L$ 对 $A$ 有正的平均因果效应（即 $L=1$ 者中 $A=1$ 的概率大于 $L=0$ 者），则在 $L\to A$ 的箭头上标正号；否则，若 $L$ 对 $A$ 有负的平均因果效应（即 $L=0$ 者中 $A=1$ 的概率大于 $L=1$ 者），则标负号。类似地，在 $L\to Y$ 的箭头上标正号或负号。若两个箭头都为正或都为负，则称混杂偏倚为正，这意味着在不校正 $L$ 时效应估计会被高估。若一个箭头为正、另一个为负，则称混杂为负，这意味着在不校正 $L$ 时效应估计会被低估。遗憾的是，在更复杂的因果图中，或变量不是二值时，这条简单规则可能失效。关于平均因果效应背景下带符号图的更详细讨论，见 VanderWeele、Hern\'an 与 Robins (2008)。

无论混杂的符号如何，另一个关键问题是偏倚的大小。不足以影响研究结论的偏倚，在实践中可以放心忽略，无论其向上还是向下。大的混杂偏倚需要强的混杂因子—处理关联与强的混杂因子—结局关联（在给定处理的条件下）。对离散的混杂因子，偏倚的大小还取决于该混杂因子的患病率（Cornfield et al. 1959; Walker 1991）。若混杂因子未知，人们只能猜测偏倚的大小。有根据的猜测可以通过敏感性分析来组织（即：就偏倚大小作出若干假设，重复进行分析），这有助于量化可以合理预期的最大偏倚。关于未测量混杂之敏感性分析的详细描述，见 Rosenbaum (2005)、Greenland (1996a)、Robins、Rotnitzky 与 Scharfstein (1999)、Greenland 与 Lash (2008)，以及 VanderWeele 与 Arah (2011)。
\end{finepoint}
\renewcommand{\FPname}{精要}

假定如上面四个例子那样，$L$、$A$、$Y$ 的数据足以识别因果效应。在这种情形下，若仅有 $A$ 与 $Y$ 的数据不足以识别（即存在结构性混杂），我们就称 $L$ 为混杂因子；若仅有 $A$、$Y$ 的数据就足以识别，我们就称 $L$ 为非混杂因子。这些定义等价于：若存在条件可交换性但不存在无条件可交换性（即结构性混杂），则称 $L$ 为混杂因子；若存在无条件可交换性，则称 $L$ 为非混杂因子。

因此，在图~7.1–7.3 中，$L$ 是混杂因子，因为 $\Pr[Y^{a}=1]$ 由标准化风险 $\sum_l \Pr[Y=1\given A=a,L=l]\Pr[L=l]$ 识别。在图~7.2 与 7.3 中，$L$ 并不是 $A$ 与 $Y$ 的共同原因，但我们仍说 $L$ 是混杂因子，因为需要它来阻断由 $A$ 与 $Y$ 的未测量共同原因 $U$ 所造成的开放后门路径。在图~7.4 中，$L$ 是非混杂因子，$\Pr[Y^{a}=1]$ 的识别公式就是条件均值 $\Pr[Y=1\given A=a]$。

\mnote{图~7.1 至 7.4 的一个非正式定义：“混杂因子就是可用于校正混杂的任何变量。”注意这一定义并不循环，因为我们此前已给出了混杂的定义。另一个非循环定义的例子：“音乐家就是演奏音乐的人。”——这句话是在我们已经定义了什么是音乐之后才说的。}

有趣的是，在图~7.4 中，给定 $L$ 的条件可交换性并不成立，因而反事实风险 $\Pr[Y^{a}=1\given L=l]$ 不等于各层内的风险 $\Pr[Y=1\given A=a,L=l]$，$L$ 各层内的条件处理效应也不可识别。此外，通过标准化 $\sum_l \Pr[Y=1\given A=a,L=l]\Pr[L=l]$ 对 $L$ 做校正，会得到 $\Pr[Y^{a}]$ 的有偏估计。这是因为对 $L$ 做校正会引入偏倚：对撞节点 $L$ 做条件化，打开了 $A$ 与 $Y$ 之间原本被该对撞节点自身阻断的后门路径（$A\leftarrow U_2\to L\leftarrow U_1\to Y$）。于是 $A$ 与 $Y$ 之间的关联将成为两种关联的混合：一种源于 $A$ 对 $Y$ 的效应，另一种源于那条开放的后门路径。关联不再等同于因果。这是我们遇到的第一个无条件可交换性成立、条件可交换性却不成立的例子：平均因果效应可以识别，但 $L$ 各水平内的条件因果效应一般不可识别。我们把条件效应中由此产生的偏倚称为选择偏倚，因为它来自对两个边际独立变量 $U_1$ 与 $U_2$ 的共同效应 $L$ 做选择（条件化），而其中一个变量与 $A$ 相关，另一个与 $Y$ 相关（见第 8 章）。\mnote{无条件效应可在不识别条件效应的情况下被识别，这一点已由 Greenland 与 Robins (1986) 以非图形方式证明。图~7.4 中的条件偏倚由 Greenland、Pearl 与 Robins (1999) 描述，并被称为 M 偏倚（Greenland 2003），因为其中涉及的变量结构——$U_2$、$L$、$U_1$——像一个横放的字母 M。}

\begin{figure}[htbp]
\centering
\begin{tikzpicture}[causal,shorten >=3pt,shorten <=1.5pt,every node/.style={font=\small}]
  \node (U1) at (0,1.2) {$U_1$};
  \node (L)  at (0,0) {$L$};
  \node (U2) at (0,-1.2) {$U_2$};
  \node (A)  at (1.7,0) {$A$};
  \node (Y)  at (3.4,0) {$Y$};
  \draw[ed] (U1) -- (L);
  \draw[ed] (U1) -- (Y);
  \draw[ed] (U2) -- (L);
  \draw[ed] (U2) -- (A);
  \draw[ed] (L) -- (A);
  \draw[ed] (A) -- (Y);
\end{tikzpicture}
\caption{后门路径与对撞路径并存}
\label{fig:7.5}
\end{figure}

图~7.5 的因果图是图~7.4 的一个变体。差别在于，图~7.5 中有一条箭头 $L\to A$。这条箭头的存在造成了一条开放的后门路径 $A\leftarrow L\leftarrow U_1\to Y$，因为 $U_1$ 是 $A$ 与 $Y$ 的共同原因，于是混杂存在。对 $L$ 做条件化会阻断那条后门路径，但同时又打开了一条以 $L$ 为对撞节点的后门路径（$A\leftarrow U_2\to L\leftarrow U_1\to Y$）。\mnote{若 $U_1$ 导致 $U_2$，或 $U_2$ 导致 $U_1$，或存在一个未测量的 $U_3$ 同时导致二者，那么 $A$ 与 $Y$ 就会存在共同原因，我们也既不会有给定 $L$ 的无条件可交换性，也不会有条件可交换性。}

因此，在图~7.5 中，偏倚难以处理：试图阻断混杂路径，却打开了选择偏倚路径。既没有无条件可交换性，也没有给定 $L$ 的条件可交换性。图~7.5 中偏倚的一个解决办法，是测量 (i) 位于 $U_1$ 与 $A$ 或 $Y$ 之间的一个变量 $L_1$，或 (ii) 位于 $U_2$ 与 $A$ 或 $L$ 之间的一个变量 $L_2$。第一种情形下，我们会有给定 $L_1$ 的条件可交换性；第二种情形下，我们会有给定 $L_2$ 与 $L$ 的条件可交换性。例如，图~7.6 包含位于 $U_1$ 与 $Y$ 之间的变量 $L_1$，以及位于 $U_2$ 与 $A$ 之间的变量 $L_2$。关于在图~7.6 中测量哪些变量会如何影响因果效应的识别，见精要~7.2。\mnote{对撞节点的定义是相对于具体路径而言的：$L$ 在路径 $A\leftarrow U_2\to L\leftarrow U_1\to Y$ 上是对撞节点，但在路径 $A\leftarrow L\leftarrow U_1\to Y$ 上不是。}

\begin{figure}[htbp]
\centering
\begin{tikzpicture}[causal,shorten >=3pt,shorten <=1.5pt,every node/.style={font=\small}]
  \node (U1) at (0,1.6) {$U_1$};
  \node (L1) at (2.0,1.3) {$L_1$};
  \node (L)  at (0,0) {$L$};
  \node (U2) at (0,-1.6) {$U_2$};
  \node (L2) at (2.0,-1.3) {$L_2$};
  \node (A)  at (3.9,0) {$A$};
  \node (Y)  at (5.7,0) {$Y$};
  \draw[ed] (U1) -- (L);
  \draw[ed] (U1) -- (L1);
  \draw[ed] (L1) -- (Y);
  \draw[ed] (L) -- (A);
  \draw[ed] (A) -- (Y);
  \draw[ed] (U2) -- (L);
  \draw[ed] (U2) -- (L2);
  \draw[ed] (L2) -- (A);
\end{tikzpicture}
\caption{在图~7.5 的结构中加入两个已测中间变量}
\label{fig:7.6}
\end{figure}

本节各因果图描绘了缺乏可交换性的两种结构性来源，二者都源于处理与结局之间存在开放的后门路径。第一种来源是处理与结局存在共同原因——它创造出一条开放的后门路径。第二种来源是对共同效应做条件化——它可能打开一条原本被阻断的后门路径。出于教学目的，我们把“混杂”一词留给前者，把“选择偏倚”留给后者。另一种从结构上定义混杂的方式，可以是“由 $A$ 与 $Y$ 之间开放的后门路径所致的偏倚”。这一替代定义与我们的定义相同，只是它把图~7.4 中因对 $L$ 做条件化而产生的偏倚称为混杂而非选择偏倚。该替代定义可以等价地表述为：混杂是“任何会被 $A$ 的随机分配所消除的系统性偏倚”。要看出这一点，请注意图~7.4 中因对 $L$ 做条件化而引入的偏倚，不可能出现在随机分配处理 $A$ 的实验中，因为随机分配保证了不存在同时为 $A$ 与 $L$ 之共同原因的未测量变量 $U_2$，于是对 $L$ 做条件化不再会打开后门路径。这两种定义之间有一个有趣的差别，如下所述。

\begin{finepoint}[条件效应与无条件效应的识别]
在任何因果图下，可识别的因果效应取决于除处理与结局之外还测量了哪些变量。以图~7.6 为例。若我们只测量 $L_2$（而不测量 $L$ 与 $L_1$），则既没有无条件可交换性，也没有给定 $L_2$ 的条件可交换性，任何因果效应都不可识别。若我们测量 $L_2$ 与 $L$，则有给定 $L_2$ 与 $L$ 的条件可交换性，但没有仅给定 $L_2$ 或仅给定 $L$ 的条件可交换性。然而，我们可以识别：

\begin{itemize}
\item $L_2$ 与 $L$ 的联合层内的条件因果效应。各反事实均值的识别公式为
\[
\E[Y\given A=a,L=l,L_2=l_2].
\]
\item 无条件因果效应。各反事实均值的识别公式为
\[
\sum_{l,l_2}\E[Y\given A=a,L=l,L_2=l_2]\,\Pr[L=l,L_2=l_2].
\]
\item $L$ 各层内的条件因果效应。各反事实均值的识别公式为
\[
\sum_{l_2}\E[Y\given A=a,L=l,L_2=l_2]\,\Pr[L_2=l_2\given L=l].
\]
\item $L_2$ 各层内的条件因果效应。各反事实均值的识别公式为
\[
\sum_{l}\E[Y\given A=a,L=l,L_2=l_2]\,\Pr[L=l\given L_2=l_2].
\]
\end{itemize}

若我们只测量 $L_1$，则有给定 $L_1$ 的条件可交换性，因而可以识别 $L_1$ 各层内的条件因果效应与无条件因果效应。若我们测量 $L_1$ 与 $L$，则还可以识别 $L_1$ 与 $L$ 的联合层内、以及 $L$ 各层内的条件因果效应。若我们测量 $L$、$L_1$ 与 $L_2$，则还可以识别这三个变量的联合层内的条件效应。
\end{finepoint}

处理与结局存在共同原因（混杂的结构定义）是关于研究总体与世界的一个实质性事实，与所选的数据分析方法无关。另一方面，把混杂定义为“会被随机化消除的任何偏倚”，则意味着混杂是否存在取决于分析方法。在图~7.4 中，若不校正 $L$，我们没有混杂；但若做校正，我们就引入了混杂。

尽管如此，在两种定义之间做选择只是口味问题，没有实际影响，因为我们关于可识别性的一切结论都只取决于条件可交换性或无条件可交换性是否成立，而不取决于我们如何定义混杂。下一章将更详细地讨论混杂与选择偏倚之别。

\section{混杂与混杂因子}

上一节描述了如何用因果图判断混杂是否存在，并在存在时确定给定的一组已测变量 $L$ 是否为混杂校正的充分集。这一程序要求事先知道因果 DAG，其中包含处理 $A$ 与结局 $Y$ 共享的一切原因——无论已测还是未测。一旦因果图已知，我们只需应用后门准则，确定需要校正哪些变量。\mnote{严格说来，研究者并不需要结构知识。他们只需要知道一个能保证条件可交换性的变量集。}

相反，处理混杂的传统方法主要基于观测到的关联，而非先验的因果知识。传统方法先把满足某些（主要是）关联条件的变量标记为混杂因子，然后规定在分析中必须对这些所谓的混杂因子做校正。当校正后的估计不同于未校正的估计时，就说存在混杂。

在传统方法下，混杂因子被定义为满足以下三个条件的变量：(1) 与处理相关；(2) 在给定处理的条件下与结局相关（“在给定处理的条件下”常被替换为“在未治疗者中”）；(3) 不位于处理与结局之间的因果路径上。然而，这种传统方法可能导致不恰当的校正。为看清原因，让我们重新审视图~7.1–7.4。

在图~7.1 中，变量 $L$ 与处理相关（因为它对 $A$ 有因果效应），在给定处理的条件下与结局相关（因为它对 $Y$ 有直接因果效应），且不位于处理与结局之间的因果路径上。在图~7.2 中，变量 $L$ 与处理相关（因为它对 $A$ 有因果效应），在给定处理的条件下与结局相关（因为它与 $Y$ 共享原因 $U$），且不位于处理与结局之间的因果路径上。在图~7.3 中，$L$ 与处理相关（它与 $A$ 共享原因 $U$），在给定处理的条件下与结局相关（它对 $Y$ 有因果效应），且不位于处理与结局之间的因果路径上。

因此，按照传统方法，在图~7.1–7.3 所表示的情形中，$L$ 是混杂因子，需要校正。这正是上一节使用后门准则所得的结论。对图~7.1–7.3，传统的、主要基于关联的方法与对因果图应用后门准则之间并无分歧。

现在再考虑图~7.4：其中没有混杂，按 7.3 节的定义 $L$ 是非混杂因子。然而，$L$ 满足传统混杂因子的判据：它与处理相关（与 $A$ 共享原因 $U_2$），在给定处理的条件下与结局相关（与 $Y$ 共享原因 $U_1$），且不位于处理与结局之间的因果路径上。于是，按照传统方法，$L$ 是应当校正的混杂因子——即便并不存在混杂！但正如上文所见，对 $L$ 做校正会因选择偏倚而使总体因果效应的估计量有偏。图~7.7 是另一个例子：传统方法导致对 $L$ 的不恰当校正，从而引入选择偏倚。

\begin{figure}[htbp]
\centering
\begin{tikzpicture}[causal,shorten >=3pt,shorten <=1.5pt,every node/.style={font=\small}]
  \node (A) at (0,0) {$A$};
  \node (L) at (1.8,0) {$L$};
  \node (Y) at (3.6,0) {$Y$};
  \node (U) at (0.5,-1.5) {$U$};
  \draw[ed] (A) -- (L);
  \draw[ed] (U) -- (L);
  \draw[ed] (U) -- (Y);
\end{tikzpicture}
\caption{传统判据导致有害校正的另一例}
\label{fig:7.7}
\end{figure}

这些例子表明，关联性的或统计性的判据不足以刻画混杂。一种几乎完全依赖统计考虑来定义混杂因子的方法，可能像图~7.4 与 7.7 所示的那样给出错误建议：即使不存在结构性混杂，也要对某个“混杂因子”做校正。为消除图~7.4 中的这一问题，传统方法的拥护者或许会把关联条件“(2) 在给定处理的条件下与结局相关”替换为结构条件“(2) 它是结局的原因”。这一修改后的混杂因子定义阻止了对图~7.4 中 $L$ 的不恰当校正，却又制造了新问题：在图~7.2 中不再把 $L$ 视为需要校正的混杂因子。见技术注记~7.2。

\begin{figure}[htbp]
\centering
\begin{tikzpicture}[causal,shorten >=3pt,shorten <=1.5pt,every node/.style={font=\small}]
  \node (U) at (0,-1.5) {$U$};
  \node (L) at (0,0) {$L$};
  \node (A) at (1.7,0) {$A$};
  \node (Y) at (3.4,0) {$Y$};
  \draw[ed] (U) -- (L);
  \draw[ed] (U) -- (A);
  \draw[ed] (U) -- (Y);
  \draw[ed] (A) -- (Y);
\end{tikzpicture}
\caption{未测量原因的替代变量}
\label{fig:7.8}
\end{figure}

传统方法会诱使研究者在校正有害时去做校正。问题之所以产生，是因为传统方法在没有充分的、关于混杂来源的因果知识时就去定义混杂因子，然后规定对这些所谓的混杂因子做校正。如果校正与未校正的估计不同，传统方法就宣称存在混杂。然而，估计的变化也可能出于混杂之外的其他原因，包括校正非混杂因子时的选择偏倚（见第 8 章），以及使用不可收缩的效应度量（见精要~4.3）。由于这些问题，基于估计变化来定义混杂的尝试早已被放弃。

相反，结构方法先明确识别混杂的来源——即处理与结局的共同原因，若它们全被测量，就足以校正混杂——然后确定一个充分的校正变量集。

结构方法清楚地表明，某个特定变量是否应纳入充分集，取决于该集已包含哪些变量。例如，在图~7.2 与 7.3 中，由于变量集 $U$ 未被测量，需要变量集 $L$ 来阻断后门路径，于是我们可以说 $L$ 中的变量是混杂因子。然而，如果 $U$ 已被测量并用于阻断后门路径，那么给定 $U$ 时 $L$ 中的变量就不再是混杂因子（另见精要~7.3）。给定一个因果 DAG，混杂是绝对的概念，而混杂因子是相对的概念。\mnote{VanderWeele 与 Shpitser (2013) 也提出了混杂因子的一个形式化定义。}

混杂的结构方法强调，从观察性数据做因果推断需要先验的因果知识。这种因果知识概括在一张因果 DAG 中，它编码了研究者关于因果网络的信念或假定。当然，无法保证研究者的因果 DAG 正确，因此有可能与研究者的信念相反，他们所选的校正变量集未能消除混杂，或引入了选择偏倚。然而，混杂的结构方法有两大优点。第一，它防止信念与行动之间的不一致。例如，如果你相信图~7.4 是真实的因果图——因而相信 $A$ 对 $Y$ 的效应没有混杂——那么无论非结构的混杂因子定义怎么说，你都不会校正变量 $L$。第二，研究者关于混杂的假定变得明确，因而可以被其他研究者明确地批评。

\begin{finepoint}[替代混杂因子]
在图~7.8 的因果 DAG 下，由于存在未测量的共同原因 $U$，$A$ 对 $Y$ 的效应存在混杂。已测变量 $L$ 是 $U$ 的代理或替代。例如，未测量的变量社会经济地位 $U$ 可能混杂体力活动 $A$ 对心血管疾病 $Y$ 之风险的效应。收入 $L$ 是常难以明确定义的变量——社会经济地位——的替代。我们应当校正变量 $L$ 吗？一方面，可以说 $L$ 不是混杂因子，因为它不在 $A$ 与 $Y$ 之间的后门路径上。另一方面，校正已测的 $L$（它与未测的 $U$ 相关）可能间接校正掉一部分由 $U$ 造成的混杂。在极端情形下，若 $L$ 与 $U$ 完全相关，那么对 $L$ 还是对 $U$ 做条件化并无差别。事实上，若 $L$ 是二值的、且是 $U$ 的无差别错分（见第 9 章）版本，则在某些弱条件下，对 $L$ 做条件化将部分阻断后门路径 $A\leftarrow U\to Y$（Greenland 1980; Ogburn and VanderWeele 2012）。因此，我们通常宁愿对 $L$ 做校正，而不是不校正。

我们把这样一些变量称为替代混杂因子：它们不在后门路径上（因而永远无法完全消除混杂），却可用于减少混杂偏倚。对抗混杂的一种可能策略，是尽可能多地测量替代混杂因子并全部校正。见第 18 章的讨论。
\end{finepoint}

\begin{techpoint}[修正传统的混杂因子定义]
图~7.4 与 7.7 给出两个图形示例，说明传统的非图形化混杂因子与混杂定义会诱使研究者去校正一个变量，而这种校正不仅多余，而且有害。传统定义之所以失败，是因为它依赖两条错误的统计判据——条件 (1) 与 (2)——和一条错误的因果判据——条件 (3)。要“修正”传统定义，需要做两件事：

\begin{enumerate}
\item 把条件 (3) 替换为：“存在变量 $L$ 与 $U$，使得在二者的联合水平内条件可交换性 $Y^{a}\indep A\given L,U$ 成立。”这一新条件比原先的条件更强，因为它实际上蕴含 $L$ 不在 $A$ 与 $Y$ 之间的因果路径上，且 $\E[Y^{a}\given L=l,U=u]$ 可由 $\E[Y\given L=l,U=u,A=a]$ 识别。
\item 把条件 (1) 与 (2) 替换为下述条件：$U$ 可分解为两个不相交的子集 $U_1$ 与 $U_2$（即 $U=U_1\cup U_2$ 且 $U_1\cap U_2$ 为空），使得 (i) 在 $L$ 的各层内 $U_1$ 与 $A$ 不相关；(ii) 在 $A$、$L$、$U_1$ 的联合层内 $U_2$ 与 $Y$ 不相关。$U_1$ 中的变量可以与 $U_2$ 中的变量相关。$U_1$ 总可以取为 $U$ 中与处理不相关的最大子集。
\end{enumerate}

若满足这两条新条件，我们就说给定 $L$ 的数据时 $U$ 是非混杂因子。这些条件由 Robins (1997a, 定理 4.3) 提出，并由 Greenland、Pearl 与 Robins (1999, 第 45–46 页) 进一步讨论（注意其中“$U=U_1\cup U_2$”这一条件被无意遗漏）。这些条件克服了图~7.4 与 7.7 中的困难，因为它们允许我们把某些变量判为非混杂因子（Robins 1997a）。例如，Greenland、Pearl 与 Robins 把这些条件应用于图~7.4，证明其中不存在混杂。
\end{techpoint}
\renewcommand{\TPname}{技术注记}

\section{单世界干预图}

可交换性译为图的语言，就是处理 $A$ 与结局 $Y$ 节点之间不存在会产生 $A$ 与 $Y$ 关联的开放路径——除那些发自 $A$ 的路径之外。第 7–9 章描述在因果图中表示缺乏可交换性的不同方式。例如，本章讨论混杂，即由于处理与结局之间存在开放的后门路径而违反可交换性。

无条件可交换性 $Y^{a}\indep A$ 与后门准则之间的等价关系显得相当神奇：反事实独立性与后门路径的不存在之间似乎没有明显联系，因为反事实并未作为变量包含在因果图中。既然图在通过 d-分离评估独立性方面如此有用，人们自然想构造把反事实也作为节点的图，从而可以直接从图上读出无条件与条件可交换性。

一种新型的图——单世界干预图（SWIG）——通过把反事实变量明确画在图上，统一了反事实途径与图形途径。SWIG 描绘的是在“所有个体都接受处理水平 $a$”的假想世界中会观测到的变量与因果关系。也就是说，SWIG 是表示由单一干预所创造的反事实世界的图。相反，标准因果图上的变量代表真实世界。于是，SWIG 可以看作一个函数，它在给定干预下变换给定的因果图。下面的例子描述这种变换。\mnote{Richardson 与 Robins (2013) 表明，SWIG 克服了此前提出的孪生因果图（Balke and Pearl 1994）的一些缺点。}

\begin{figure}[htbp]
\centering
\begin{tikzpicture}[causal,shorten >=3pt,shorten <=1.5pt,every node/.style={font=\small}]
  \node (U) at (0,-1.4) {$U$};
  \node (L) at (0,0) {$L$};
  \node (A) at (1.9,0) {$A\,|\,a$};
  \node (Y) at (3.7,0) {$Y^{a}$};
  \draw[ed] (U) -- (L);
  \draw[ed] (U) -- (Y);
  \draw[ed] (L) -- (A);
  \draw[ed] (A) -- (Y);
\end{tikzpicture}
\caption{由图~7.2 得到的单世界干预图}
\label{fig:7.9}
\end{figure}

设图~7.2 的因果图代表观测到的研究数据。图~7.9 的 SWIG 是图~7.2 的一种变换，表示所有个体都接受了“把其处理设为固定值 $a$”的干预的世界。

在 SWIG 中，处理节点被分成左右两半，一旦分开，二者就被视为各自独立的节点（变量）。右半编码干预之下的处理取值 $a$，并继承原因果 DAG 中所有从 $A$ 出发的箭头。左半编码在无干预时会观测到的处理 $A$ 的取值，即处理的自然取值。它继承因果 DAG 上所有指向 $A$ 的节点，因为在被干预的（反事实）世界中，其因果输入与真实世界相同。注意 $A$ 没有指向 $a$ 的箭头，因为取值 $a$ 对所有人都相同，即在被干预的世界中是一个常数。

我们假定处理的自然取值 $A$ 定义良好，尽管在干预 $a$ 之下我们通常无法测量它。不过在有些情形中，$A$ 也许可以测量：近来的实验提示，脑电图记录能够在个体意识到自己的决定之前最多 1/2 秒，探测到个体将要做出的选择。若果真如此，就可以通过脑电图实际测量 $A$，同时还留下 1/2 秒去实施干预、给予处理 $a$。\mnote{在单干预世界中，$a$ 是常数，不会影响其他变量。然而在绘制 SWIG 时，我们从 $a$ 引出箭头，作为一种便捷方式，以追踪原 DAG 中被 $A$ 直接影响的变量。}

\begin{figure}[htbp]
\centering
\begin{tikzpicture}[causal,shorten >=3pt,shorten <=1.5pt,every node/.style={font=\small}]
  \node (U1) at (0,1.2) {$U_1$};
  \node (L)  at (0,0) {$L$};
  \node (U2) at (0,-1.2) {$U_2$};
  \node (A)  at (2.0,0) {$A\,|\,a$};
  \node (Y)  at (3.9,0) {$Y^{a}$};
  \draw[ed] (U1) -- (L);
  \draw[ed] (U1) -- (Y);
  \draw[ed] (U2) -- (L);
  \draw[ed] (U2) -- (A);
  \draw[ed] (A) -- (Y);
\end{tikzpicture}
\caption{由图~7.4 得到的单世界干预图}
\label{fig:7.10}
\end{figure}

在 SWIG 中，结局是 $Y^{a}$，即被干预世界中的 $Y$ 取值。由于其余变量在时间上先于 $A$，它们不受干预影响，因而取与观测世界相同的取值，即不被标记为反事实变量。事实上，任何非 $A$ 之子孙的变量都不必标记为反事实，因为在忠实性假定（我们采用这一假定）下，处理对任何个体的非子孙变量都没有因果效应。在我们的因果模型下，条件可交换性 $Y^{a}\indep A\given L$ 成立，因为对 $L$ 做条件化后，$Y^{a}$ 与 $A$ 之间的所有路径都被阻断，即给定 $L$ 时 $Y^{a}$ 与 $A$ 是 d-分离的。\mnote{在 FFRCISTG 模型下可以证明，d-分离也蕴含 SWIG 上的统计独立性。}

现在考虑图~7.4 的因果图与图~7.10 的 SWIG。边际可交换性 $Y^{a}\indep A$ 成立，因为在 SWIG 上，$Y^{a}$ 与 $A$ 之间的所有路径都被阻断（未对 $L$ 做条件化）。相反，条件可交换性 $Y^{a}\indep A\given L$ 不成立，因为在 SWIG 上，当对撞节点 $L$ 被条件化时，路径 $Y^{a}\leftarrow U_1\to L\leftarrow U_2\to A$ 是开放的。这就是为什么边际的 $A$–$Y$ 关联是因果的，而给定 $L$ 的条件 $A$–$Y$ 关联不是；也因此，任何校正 $L$ 的方法都会导致偏倚。这些例子显示了 SWIG 如何统一反事实途径与图形途径。事实上不难看出，在 SWIG 上，$Y^{a}$ 与 $A$ 在给定 $L$ 时 d-分离，当且仅当 $L$ 是 $A$ 的非子孙变量且阻断了从 $A$ 到 $Y$ 的所有后门路径（另见精要~7.4）。

\section{混杂的调整}

在没有随机化的情况下，因果推断依赖于一个无法检验的假定：我们已测量了一组变量 $L$，它是混杂校正的充分集，即一组 $A$ 的非子孙变量，包含足够多的变量以阻断从 $A$ 到 $Y$ 的所有后门路径。在给定 $L$ 的条件可交换性这一假定下，可以用标准化与 IP 加权计算总体中的平均因果效应。但正如 4.6 节所讨论的，标准化与 IP 加权并非观察性研究中校正混杂的唯一方法。校正混杂因子 $L$ 的方法可分为两大类：

\begin{itemize}
\item G 方法：标准化、IP 加权与 g 估计。这些方法（“g”代表“广义”）利用给定 $L$ 的条件可交换性，估计 $A$ 对 $Y$ 在整个总体或总体任一子集中的因果效应。在我们的心脏移植研究中，我们在 2.4 节（标准化）与 2.5 节（IP 加权）用 G 方法来校正疾病严重程度 $L$ 造成的混杂。第二部分描述 G 方法的模型化扩展：参数 g 公式（标准化）、边际结构模型的 IP 加权，以及嵌套结构模型的 g 估计。
\item 基于分层的传统校正方法：分层（包括限制）与匹配。这些方法利用给定 $L$ 的条件可交换性，估计由 $L$ 所定义的各子集中 $A$ 与 $Y$ 之间的关联。在我们的心脏移植研究中，我们在 4.4 节（分层）与 4.5 节（匹配）用基于分层的方法校正疾病严重程度 $L$ 造成的混杂。第二部分描述传统分层的模型化扩展：结局回归。
\end{itemize}

标准化与 IP 加权模拟的是：若涉及已测变量 $L$ 的后门路径不存在，总体中的 $A$–$Y$ 关联会是什么样。例如，IP 加权通过构造一个伪总体来实现这一点：其中处理 $A$ 与已测混杂因子 $L$ 独立，即“删去”从 $L$ 到 $A$ 的箭头。相反，基于分层的传统方法不删去从 $L$ 到 $A$ 的箭头，而是在观测总体的一个子集中计算条件效应，这相当于添加一个选择框。在第三部分（聚焦于时变处理），我们将描述为什么在使用标准化或 IP 加权时“删去”箭头 $L\to A$ 有优势，以及为什么 g 估计是唯一普遍有效的基于分层的方法。传统基于分层方法的偏倚将在第 20 章描述。在时变处理（因而时变混杂因子）的情形中，G 方法是校正混杂的首选，因为传统基于分层的方法可能导致选择偏倚。\mnote{分层与匹配的一种常见变体，是把每个个体的变量 $L$ 替换为该个体接受治疗的估计概率 $\Pr[A=1\given L]$，即倾向性评分（Rosenbaum and Rubin 1983）。见第 15 章。}

上述所有方法都要求给定 $L$ 的条件可交换性。然而，混杂有时可以用不要求条件可交换性的方法来处理。这类方法的例子有双重差分（技术注记~7.3）、工具变量估计（第 16 章）、近端推断（技术注记~7.3）、前门准则（技术注记~7.4）等。遗憾的是，这些方法要求替代性的假定，而这些假定与条件可交换性一样无法验证。因此在实践中，所得效应估计的有效性并无保证。校正方法的选择，取决于在特定情境中，哪一组无法验证的假定——条件可交换性还是替代条件——被认为更可能成立。

在许多观察性研究中，达到条件可交换性可能是不现实的目标；但正如 3.2 节所讨论的，关于因果结构的专家知识可以用来尽可能接近这一目标。因此，在观察性研究中，研究者测量许多变量 $L$（它们是处理的非子孙），试图确保受处理者与未处理者条件可交换。所希望的是：即便可能存在共同原因（混杂），已测变量 $L$ 也足以阻断所有后门路径（无未测量混杂）。然而，无法保证这一尝试会成功，这使从观察性数据做因果推断成为一件有风险的事。

此外，专家知识可用于避免校正那些可能引入偏倚的变量。至少，研究者一般应避免校正受处理或结局影响的变量。当然，思虑周密、知识渊博的研究者可能认为两种或更多的因果结构同样可信，而它们可能导致关于混杂与混杂因子的不同结论。在这种情况下，他们会做多种分析，并明确陈述每种分析有效所需的因果结构假定。遗憾的是，人们永远无法确定所考虑的因果结构集合包含了真实的那一个；这种不确定性在观察性数据中是不可避免的。\mnote{运用因果结构的专家知识校正混杂的一个实际例子，由 Hern\'an 等 (2002) 描述。}

观察性研究中始终存在的混杂威胁，有一个科学上的后果。设想你做了一项观察性研究，以量化心脏移植 $A$ 对死亡 $Y$ 的效应。你尽了最大努力（例如咨询领域专家）去识别并测量混杂因子 $L$，并假定校正 $L$ 后不存在未测量混杂。一位批评者说：“这项观察性研究的推断可能因潜在混杂而错误。”这位批评者所做的不是科学陈述，而是逻辑陈述。既然任何观察性研究的结果都可能受到混杂，那么你的研究结果可能受到混杂，显然是成立的。如果批评者的本意是提供关于你这项具体研究之缺陷的证据，那么他失败了。他的批评没有提供信息，因为他只是重述了观察性研究的一个特征，而这个特征在你和批评者在研究开始之前就已知道。

要恰当地批评你的研究，批评者需要投入一场科学对话。例如，他可以引用与你的发现相矛盾的实验或观察证据，或者他可以说：“这项观察性研究的推断可能因吸烟的潜在混杂而错误，吸烟是一个共同原因，经由它可能留下一条开放的后门路径。”后一种选择为你关于“无未测量混杂”的假定提供了一个可检验的挑战。举证责任又落到你身上。你的下一步是尝试校正吸烟；或者，如果无法获得吸烟数据，就做敏感性分析，考察吸烟可能诱导的偏倚。

尽管上述讨论仅限于混杂所致的偏倚，要从观察性数据做有效的因果推断，还需要不存在选择偏倚与测量偏倚。但与混杂不同，这些其他偏倚在随机化实验与观察性研究中都可能出现。在探讨了本章的混杂之后，下一章将呈现受处理者与未受处理者之间缺乏可交换性的另一个潜在来源：个体被选入分析。

\begin{finepoint}[混杂因子不能是处理的子孙，却可以位于处理的未来]
考虑图~7.11 的因果 DAG。$L$ 是处理 $A$ 的子孙，它阻断了从 $A$ 到 $Y$ 的所有后门路径。与图~7.4 和 7.7 不同，对 $L$ 做条件化不会造成选择偏倚，因为没有打开任何对撞路径。相反，由于 $A$ 对 $Y$ 的因果效应完全通过中间变量 $L$，对 $L$ 做条件化会完全阻断这条通路。这个例子表明，当 $L$ 是 $A$ 的子孙时，校正一个阻断了所有后门路径的变量 $L$ 并不能消除偏倚。

既然条件可交换性 $Y^{a}\indep A\given L$ 蕴含对 $L$ 的校正能消除一切偏倚，那么在本例中必定是条件可交换性不成立，且平均因果效应 $\E[Y^{a=1}]-\E[Y^{a=0}]$ 不可识别。这一失败可以通过分析图~7.12 的 SWIG 加以验证，它描绘了把 $A$ 设为取值 $a$ 的反事实世界。在这个世界里，事实性变量 $L$ 被反事实变量 $L^{a}$ 取代，即若所有个体都接受处理取值 $a$ 时本会观测到的 $L$ 取值。由于 $L^{a}$ 阻断了从 $Y^{a}$ 到 $A$ 的所有路径，我们断定 $Y^{a}\indep A\given L^{a}$ 成立；但我们不能由此断定条件可交换性 $Y^{a}\indep A\given L$ 成立，因为 $L$ 根本不在图上。（在 FFRCISTG 下，任何无法从 SWIG 上读出的独立性都不能被假定成立。）因此，我们无法确保平均因果效应 $\E[Y^{a=1}]-\E[Y^{a=0}]$ 可由 $(L,A,Y)$ 的数据识别。

问题之所以产生，是因为 $L$ 是 $A$ 的子孙，而不是因为 $L$ 位于 $A$ 的未来。如果在图~7.11 中从 $A$ 到 $L$ 的箭头不存在，那么 $L$ 就是 $A$ 的非子孙变量，且阻断所有后门路径。类似地，在图~7.12 的 SWIG 上，由于 $A$ 不再是 $L$ 的原因，我们可以把 $L^{a}$ 换回 $L$（注意此时 $Y^{a}$ 与 $A$ 被 $L$ 所 d-分离）。因此，即便 $L$ 仍位于 $A$ 的未来，校正 $L$ 也会消除一切偏倚。关键在于因果图的拓扑结构（哪个变量导致哪个变量），而不在于节点的时间顺序。Rosenbaum (1984) 与 Robins (1986, 第 11 节) 对用时点后变量控制混杂做了非图形化的讨论。
\end{finepoint}

\begin{figure}[htbp]
\centering
\begin{tikzpicture}[causal,shorten >=3pt,shorten <=1.5pt,every node/.style={font=\small}]
  \node (A) at (0,0) {$A$};
  \node (L) at (1.7,0) {$L$};
  \node (Y) at (3.4,0) {$Y$};
  \node (U) at (0.5,-1.4) {$U$};
  \draw[ed] (A) -- (L);
  \draw[ed] (L) -- (Y);
  \draw[ed] (U) -- (A);
  \draw[ed] (U) -- (L);
\end{tikzpicture}
\caption{处理的中介变量}
\label{fig:7.11}
\end{figure}

\begin{figure}[htbp]
\centering
\begin{tikzpicture}[causal,shorten >=3pt,shorten <=1.5pt,every node/.style={font=\small}]
  \node (U) at (0,-1.4) {$U$};
  \node (A) at (0,0) {$A\,|\,a$};
  \node (L) at (2.0,0) {$L^{a}$};
  \node (Y) at (3.8,0) {$Y^{a}$};
  \draw[ed] (U) -- (A);
  \draw[ed] (U) -- (L);
  \draw[ed] (A) -- (L);
  \draw[ed] (L) -- (Y);
\end{tikzpicture}
\caption{图~7.11 的单世界干预图}
\label{fig:7.12}
\end{figure}

\begin{techpoint}[双重差分与负结局对照]
设我们想计算阿司匹林 $A$（1：是；0：否）对血压 $Y$ 的平均因果效应，但存在 $A$ 与 $Y$ 的未测量共同原因 $U$，例如心脏病史。那么我们就无法通过标准化或 IP 加权计算该效应，因为存在未测量混杂。但另有一种方法，在一定条件下可以校正未测量混杂：使用负结局对照（也称“安慰剂检验”）。

再设对总体中每个个体，我们还测量了治疗可得之前那一刻的结局取值。我们把这一处理前结局 $C$ 称为负结局对照（也称阴性对照结局）。如图~7.13 所示，$U$ 是 $Y$ 与 $C$ 的共同原因，而处理 $A$ 显然不是处理前 $C$ 的原因。现在，尽管已知 $A$ 对 $C$ 的因果效应为零，对比 $\E[C\given A=1]-\E[C\given A=0]$ 却不为零，因为存在 $U$ 造成的混杂。事实上，$\E[C\given A=1]-\E[C\given A=0]$ 在加法尺度上度量了 $A$ 对 $C$ 之效应的混杂大小。若 $A$ 对负结局对照 $C$ 之效应的加法混杂大小，与 $A$ 对真实结局 $Y$ 之效应的加法混杂大小相同，那么我们就可以计算 $A$ 对 $Y$ 在处理者中的效应。具体而言，在加法等量混杂假定
\[
\E[Y^{1}-Y^{0}\given A=1]=\E[C\given A=1]-\E[C\given A=0]
\]
下，该效应为
\[
\begin{aligned}
\E[Y^{1}-Y^{0}\given A=1]={}&\bigl(\E[Y\given A=1]-\E[Y\given A=0]\bigr)\\
&-\bigl(\E[C\given A=1]-\E[C\given A=0]\bigr).
\end{aligned}
\]
也就是说，处理者中的效应等于处理 $A$ 与结局 $Y$ 之间的关联（它是因果效应与混杂的混合）减去由 $A$ 与 $C$ 之间的关联所度量的混杂。注意，图~7.13 中从 $C$ 到 $Y$ 的直接箭头对于 $C$ 作为负结局对照并非必需。

这种校正混杂的方法称为双重差分（Card 1990; Meyer 1995; Angrist and Krueger 1999）。在实践中，该方法常与使用参数或半参数方法对已测协变量的校正结合使用（Abadie 2005）。然而，双重差分是负结局对照的一种略显局限的做法（Sofer et al. 2016）：它要求同时测量处理前与处理后的结局（或至少要求真实结局 $Y$ 与负结局对照 $C$ 在同一尺度上测量），并要求加法等量混杂。Sofer 等 (2016) 描述了更一般的方法，允许 $Y$ 与 $C$ 处于不同尺度，依赖较弱版本的等量混杂，并纳入对已测协变量的校正。关于用负结局对照检测混杂的一般介绍，见 Lipsitch et al. (2010) 与 Flanders et al. (2011)。

令人惊讶的是，当同时拥有负结局对照 $C$ 与负处理对照 $Z$ 时，即便存在未测量混杂因子 $U$，在附加假定下因果效应也可以非参数识别。事实上，若 $U$、$C$、$Z$ 都是离散的，且 $C$ 与 $Z$ 的水平数至少与 $U$ 一样多，则 $A$ 对 $Y$ 的因果效应相当普遍地可以被识别（Miao et al. 2018）。这种识别途径称为近端因果推断（Cui et al. 2024）。图~7.15 是一个例子，其中 $C$ 是负结局对照，$Z$ 是负处理对照。
\end{techpoint}

\begin{figure}[htbp]
\centering
\begin{tikzpicture}[causal,shorten >=3pt,shorten <=1.5pt,every node/.style={font=\small}]
  \node (C) at (0,0.5) {$C$};
  \node (A) at (2.3,0.5) {$A$};
  \node (Y) at (4.2,0.5) {$Y$};
  \node (U) at (0,-1.1) {$U$};
  \draw[ed] (U) -- (C);
  \draw[ed] (U) -- (A);
  \draw[ed] (U) -- (Y);
  \draw[ed] (A) -- (Y);
  \draw[ed] (C) to[out=50,in=130] (Y);
\end{tikzpicture}
\caption{负结局对照}
\label{fig:7.13}
\end{figure}

\begin{figure}[htbp]
\centering
\begin{tikzpicture}[causal,shorten >=3pt,shorten <=1.5pt,every node/.style={font=\small}]
  \node (U) at (0,-1.3) {$U$};
  \node (A) at (1.8,0) {$A$};
  \node (M) at (3.2,0) {$M$};
  \node (Y) at (4.6,0) {$Y$};
  \draw[ed] (U) -- (A);
  \draw[ed] (U) -- (Y);
  \draw[ed] (A) -- (M);
  \draw[ed] (M) -- (Y);
\end{tikzpicture}
\caption{前门结构}
\label{fig:7.14}
\end{figure}

\begin{figure}[htbp]
\centering
\begin{tikzpicture}[causal,shorten >=3pt,shorten <=1.5pt,every node/.style={font=\small}]
  \node (C) at (0,0.6) {$C$};
  \node (Z) at (1.9,0.6) {$Z$};
  \node (A) at (3.5,0.6) {$A$};
  \node (Y) at (5.3,0.6) {$Y$};
  \node (U) at (0,-1.1) {$U$};
  \draw[ed] (U) -- (C);
  \draw[ed] (U) -- (Z);
  \draw[ed] (U) -- (A);
  \draw[ed] (U) -- (Y);
  \draw[ed] (Z) -- (A);
  \draw[ed] (A) -- (Y);
  \draw[ed] (C) to[out=50,in=130] (Y);
\end{tikzpicture}
\caption{负结局对照与负处理对照}
\label{fig:7.15}
\end{figure}

\begin{techpoint}[前门准则]
图~7.14 的因果图描绘了这样一种情形：处理 $A$ 与二值结局 $Y$ 共享一个未测量的原因 $U$，且存在一个变量 $M$，它完全中介 $A$ 对 $Y$ 的效应，并与 $A$ 或 $Y$ 都不共享未测量原因。在这种因果结构下，数据分析者无法直接用标准化（也无法用 IP 加权）计算反事实风险 $\Pr[Y^{a=1}=1]$ 与 $\Pr[Y^{a=0}=1]$，因为阻断 $A$ 与 $Y$ 之间后门路径所必需的变量 $U$ 不可得。因此，用前几章所述的方法无法识别 $A$ 对 $Y$ 的平均因果效应。然而，Pearl (1995) 表明，$\Pr[Y^{a}=1]$ 可由所谓的前门公式识别：
\[
\Pr[Y^{a}=1]=\sum_m \Pr[M=m\given A=a]\sum_{a'}\Pr[Y=1\given M=m,A=a']\,\Pr[A=a'].
\]

Pearl 把这一识别公式称为前门校正，因为它依赖于一条从 $A$ 到 $Y$ 的路径的存在；与后门路径相反，这条路径经过 $A$ 的子孙 $M$，而 $M$ 完全中介了 $A$ 对 $Y$ 的效应。Pearl 常用“后门公式”一词来指我们称为标准化或（点处理）g 公式的识别公式（Robins 1986）。下面给出前门识别公式的证明。

注意
\[
\Pr[Y^{a}=1]=\sum_m \Pr[M^{a}=m]\,\Pr[Y^{a}=1\given M^{a}=m],
\]
且在图~7.14 下，$\Pr[M^{a}=m]=\Pr[M=m\given A=a]$，因为 $A$ 对 $M$ 的效应没有混杂（即 $A\indep M^{a}$）；又 $\Pr[Y^{a}=1\given M^{a}=m]=\Pr[Y^{m}=1]$，因为 (i) 当 $M^{a}=m$ 时 $Y^{a}=Y^{m}$（在图~7.14 中 $A$ 只通过 $M$ 影响 $Y$），(ii) 在把 $M$ 设为 $m$、$A$ 设为 $a$ 的联合干预下，由 SWIG 上的 d-分离有 $Y^{m}\indep M^{a}$。最后，在只对 $M$ 做干预的 SWIG 上，由条件可交换性 $Y^{m}\indep M\given A$，得
\[
\Pr[Y^{m}=1]=\sum_{a'}\Pr[Y=1\given M=m,A=a']\,\Pr[A=a'].
\]

上述证明要求在对 $M$ 的干预下反事实结局 $Y^{m}$ 定义良好。在技术注记~21.11 与 21.12 中，我们给出前门公式的替代证明，它们不要求这一条件。
\end{techpoint}
